% CUMCM 国赛数学建模论文骨架（中文，XeLaTeX 编译：latexmk -xelatex main.tex）
\documentclass[12pt]{article}
\usepackage{ctex}                 % 中文支持
\usepackage{geometry}
\geometry{a4paper, margin=2.5cm}
\usepackage{amsmath, amssymb}
\usepackage{graphicx}
\usepackage{booktabs}             % 三线表
\usepackage{caption}
\usepackage{float}
\usepackage[numbers,sort&compress]{natbib}
\usepackage{hyperref}

\title{\textbf{基于……的……模型}}
\author{}
\date{}

\begin{document}

% ===== 摘要（评审第一关：方法 + 关键结果 + 结论） =====
\begin{center}\Large\textbf{摘要}\end{center}
\noindent 针对问题一，建立了……模型，采用……方法，求得……（关键数据）。
针对问题二，……。针对问题三，……。本文模型的优点在于……。

\noindent\textbf{关键词：} 关键词1；关键词2；关键词3
\newpage

\section{问题重述}
\subsection{问题背景}
\subsection{问题提出}

\section{问题分析}

\section{模型假设}
\begin{enumerate}
  \item 假设……，理由：……。
  \item 假设……，理由：……。
\end{enumerate}

\section{符号说明}
\begin{table}[H]\centering
\begin{tabular}{cll}
\toprule
符号 & 含义 & 单位 \\
\midrule
$x$ & …… & …… \\
\bottomrule
\end{tabular}
\end{table}

\section{模型的建立与求解}
\subsection{问题一模型}
\begin{equation}
  \min\ f(x)\quad \text{s.t.}\ g(x)\le 0
\end{equation}
% \begin{figure}[H]\centering\includegraphics[width=.7\textwidth]{figure.pdf}
% \caption{结果图}\end{figure}

\section{模型检验与灵敏度分析}

\section{模型评价}
\subsection{优点}
\subsection{缺点与改进}

\bibliographystyle{gbt7714-numerical} % 或 plainnat
\bibliography{refs}

\appendix
\section{支撑代码}
% \lstinputlisting{solve.py}

\end{document}
